\chapter{From Linear Algebra to Many-Body Physics}
\label{chap:mbphys}
%=============================================================================

Chapter~\ref{chap:linalg} developed a language and a set of algorithms.  This
chapter puts them to work.  We shall see that the central object of fermionic
many-body theory, the Slater determinant, is a determinant in exactly the
sense of Section~\ref{sec:determinants}; that changing the single-particle
basis is a unitary transformation of the kind studied in
Section~\ref{sec:unitary}; that the Hartree-Fock energy is minimised by
solving a symmetric eigenvalue problem with the machinery of
Sections~\ref{sec:householder} and \ref{sec:qlalgorithm}; and that the
obstacle which makes the whole subject difficult -- the exponential growth of
the Hilbert space -- is the tensor product of Section~\ref{sec:tensor}.

Nothing in this chapter is new mathematics.  What is new is the physics that
the mathematics is asked to carry.

%----------------------------------------------------------------------
\section{The many-body Hamiltonian}
\label{sec:hamiltonian}

Throughout these lectures we assume that the interacting part of the
Hamiltonian can be approximated by a two-body interaction, so the full
$N$-particle Hamiltonian takes the form
\begin{equation}
  \hat{H} = \hat{H}_0 + \hat{H}_I
           = \sum_{i=1}^{N}\hat{h}_0(x_i)
           + \sum_{i<j}^{N}\hat{v}(r_{ij}),
  \label{eq:H}
\end{equation}
where $x_i$ denotes the coordinates -- position and spin -- of particle $i$.
The one-body operator splits into a kinetic and a potential part,
\begin{equation}
  \hat{h}_0(x_i) = \hat{t}(x_i) + \hat{u}_{\mathrm{ext}}(x_i),
  \label{eq:2-onebody}
\end{equation}
where $\hat{u}_{\mathrm{ext}}$ is typically a harmonic oscillator potential,
the Coulomb attraction of a nucleus, or the self-consistent Hartree-Fock
potential derived later in this chapter.

We shall assume that the single-particle functions $\psi_{\alpha}$ are
eigenfunctions of $\hat{h}_0$,
\begin{equation}
  \hat{h}_0(x)\,\psi_{\alpha}(x)
   = \left(\hat{t}(x)+\hat{u}_{\mathrm{ext}}(x)\right)\psi_{\alpha}(x)
   = \varepsilon_{\alpha}\,\psi_{\alpha}(x),
  \label{eq:2-spequation}
\end{equation}
with $\varepsilon_{\alpha}$ the \emph{unperturbed} or non-interacting
single-particle energies.  In the absence of $\hat{H}_I$ the total energy
would simply be the sum of the occupied $\varepsilon_{\alpha}$.  Everything
that makes many-body physics hard is contained in the correction to that
statement.

Equation~(\ref{eq:2-spequation}) is an eigenvalue problem, and in practice we
solve it exactly as in Section~\ref{sec:schrodiag}: discretise on a grid, or
expand in a finite basis, and diagonalise the resulting matrix with the
algorithms of Sections~\ref{sec:householder} and~\ref{sec:qlalgorithm}.  The
harmonic oscillator eigenfunctions computed there are the single-particle
basis used in the numerical examples of this chapter.

Finally, the Hamiltonian~(\ref{eq:H}) is \emph{symmetric} under the
interchange of any two particles.  Writing $\hat{P}_{ij}$ for the operator
that permutes particles $i$ and $j$,
\begin{equation}
  [\hat{H},\hat{P}_{ij}] = 0 ,
  \label{eq:2-commuteP}
\end{equation}
so the eigenfunctions of $\hat{H}$ can be chosen to be eigenfunctions of
$\hat{P}_{ij}$ as well.

%----------------------------------------------------------------------
\section{The Pauli exclusion principle and antisymmetry}
\label{sec:pauli_exc}

Let $\Psi_{\lambda}$ be such a simultaneous eigenfunction.  Then
\begin{equation}
  \hat{P}_{ij}\Psi_{\lambda}(x_1,\dots,x_i,\dots,x_j,\dots,x_N)
   = \beta\,\Psi_{\lambda}(x_1,\dots,x_i,\dots,x_j,\dots,x_N),
  \label{eq:2-permeigen}
\end{equation}
and since applying the permutation twice returns the original function,
$\beta^2=1$ and $\beta=\pm1$.  For bosons $\beta=+1$; for fermions, which are
our concern here, the Pauli principle requires
\begin{equation}
  \beta = -1 ,
  \label{eq:2-beta}
\end{equation}
so the wave function changes sign under every transposition.

An arbitrary product of single-particle functions has no such symmetry, and
must be projected onto the antisymmetric subspace.  The projector is the
\emph{antisymmetrisation operator}
\begin{equation}
  \hat{A} = \frac{1}{N!}\sum_{p}(-)^{p}\hat{P},
  \label{eq:2-antisym}
\end{equation}
where the sum runs over all $N!$ permutations $\hat{P}$ and $p$ is the parity
of each, that is the number of transpositions required to build it.  Two
properties of $\hat{A}$ will do all the work in Section~\ref{sec:expval}.
First, since $\hat{H}_0$ and $\hat{H}_I$ are invariant under every
permutation,
\begin{equation}
  [\hat{H}_0,\hat{A}] = [\hat{H}_I,\hat{A}] = 0 .
  \label{eq:2-commuteA}
\end{equation}
Second, $\hat{A}$ is idempotent,
\begin{equation}
  \hat{A}^2 = \hat{A},
  \label{eq:Asq}
\end{equation}
because permuting an already antisymmetrised function reproduces it up to the
sign that the prefactor $(-)^p$ removes.  A projector, in the sense of
Section~\ref{sec:compbasis}: $\hat{A}$ projects the full $N$-particle space
onto its antisymmetric subspace, exactly as $\bm{P}=\ketbra{0}{0}$ projects a
qubit onto a computational basis state.

%----------------------------------------------------------------------
\section{Determinants, and how they are actually computed}
\label{sec:determinants}

For an $n\times n$ matrix $\bm{A}$, the determinant is defined by summing over
all $n!$ permutations of the column indices,
\begin{equation}
  \Det{A} = |\bm{A}|
   = \sum_{i=1}^{n!}(-1)^{p_i}\hat{P}_i\,a_{11}a_{22}\cdots a_{nn},
  \label{eq:2-detdef}
\end{equation}
where $\hat{P}_i$ permutes the column indices and $p_i$ is the number of
transpositions needed to restore the natural ordering $a_{11}a_{22}\cdots
a_{nn}$.  For $n=2$ the sum has two terms,
\begin{equation}
  \begin{vmatrix}a_{11}&a_{12}\\a_{21}&a_{22}\end{vmatrix}
   = (-1)^0 a_{11}a_{22} + (-1)^1 a_{12}a_{21},
  \label{eq:2-det2}
\end{equation}
the second obtained by interchanging the column indices $1$ and $2$.

The resemblance between Eq.~(\ref{eq:2-detdef}) and the antisymmetriser
(\ref{eq:2-antisym}) is not an accident, and it is the reason determinants
appear in fermionic physics at all: both are the signed sum over the symmetric
group, and the determinant is simply what the antisymmetriser produces when
applied to a product of single-particle functions.

\paragraph{Two properties we shall need.}
The determinant is linear in each column separately,
\begin{equation}
  \left|\begin{array}{cccc}
    a_{11} & \cdots & \displaystyle\sum_{k=1}^{n}c_k b_{1k} & \cdots \\[4pt]
    a_{21} & \cdots & \displaystyle\sum_{k=1}^{n}c_k b_{2k} & \cdots \\
    \vdots & & \vdots & \\
    a_{n1} & \cdots & \displaystyle\sum_{k=1}^{n}c_k b_{nk} & \cdots
  \end{array}\right|
  =\sum_{k=1}^{n}c_k
  \left|\begin{array}{cccc}
    a_{11} & \cdots & b_{1k} & \cdots \\
    a_{21} & \cdots & b_{2k} & \cdots \\
    \vdots & & \vdots & \\
    a_{n1} & \cdots & b_{nk} & \cdots
  \end{array}\right| ,
  \label{eq:2-detlinear}
\end{equation}
and if every column of $\bm{A}$ is a linear combination of the columns of
$\bm{B}$ with coefficient matrix $\bm{C}$, that is $a_{ij}=\sum_k b_{ik}c_{kj}$
or $\bm{A}=\bm{B}\bm{C}$, then
\begin{equation}
  \Det{A} = \Det{B}\cdot\Det{C} .
  \label{eq:2-detproduct}
\end{equation}
Equation~(\ref{eq:2-detproduct}) follows by applying
Eq.~(\ref{eq:2-detlinear}) to each column in turn; the reader is invited to
verify it for $n=2$.  It is the single most important algebraic fact in this
chapter, and Section~\ref{sec:basischange} is nothing but its physical
interpretation.

\paragraph{A warning about Eq.~(\ref{eq:2-detdef}).}
The definition is a definition, not an algorithm.  Evaluating a determinant by
summing over permutations costs $\bigO(n!)$ operations, which for $n=20$ is
$2\times10^{18}$ and for $n=25$ exceeds the age of the universe on any
machine.  The LU decomposition of Section~\ref{sec:lu} computes the same
number in $\bigO(n^3)$ operations as the product of the diagonal elements of
$\bm{U}$, Eq.~(\ref{eq:1-ludet}).  Table~\ref{tab:detcost} makes the
comparison.  Every Slater determinant evaluated in a Monte Carlo or
configuration-interaction code is computed this way, never from
Eq.~(\ref{eq:2-detdef}).

\begin{table}[h]
\centering
\renewcommand{\arraystretch}{1.25}
\setlength{\tabcolsep}{12pt}
\begin{tabular}{rll}
\toprule
$n$ & \textbf{From Eq.~(\ref{eq:2-detdef}), $n!$} &
\textbf{From LU, $\sim\tfrac{2}{3}n^3$}\\
\midrule
 5 & $1.2\times10^{2}$  & $8.3\times10^{1}$\\
10 & $3.6\times10^{6}$  & $6.7\times10^{2}$\\
20 & $2.4\times10^{18}$ & $5.3\times10^{3}$\\
50 & $3.0\times10^{64}$ & $8.3\times10^{4}$\\
\bottomrule
\end{tabular}
\caption{Operation counts for evaluating an $n\times n$ determinant from the
definition and from the LU decomposition of Section~\ref{sec:lu}.}
\label{tab:detcost}
\end{table}

%----------------------------------------------------------------------
\section{The Slater determinant}
\label{sec:slaterdet}

We now have everything needed to write down an antisymmetric $N$-particle
wave function.  Given $N$ single-particle orbitals
$\psi_{\alpha},\psi_{\beta},\dots,\psi_{\sigma}$, the \emph{Slater
determinant} is
\begin{equation}
  \Phi(x_1,\dots,x_N;\alpha,\beta,\dots,\sigma)
  =\frac{1}{\sqrt{N!}}
  \begin{vmatrix}
    \psi_\alpha(x_1) & \psi_\alpha(x_2) & \cdots & \psi_\alpha(x_N)\\
    \psi_\beta(x_1)  & \psi_\beta(x_2)  & \cdots & \psi_\beta(x_N)\\
    \vdots & \vdots & \ddots & \vdots\\
    \psi_\sigma(x_1) & \psi_\sigma(x_2) & \cdots & \psi_\sigma(x_N)
  \end{vmatrix},
  \label{eq:SlatDet}
\end{equation}
where $\alpha,\beta,\dots,\sigma$ are the quantum numbers labelling the
occupied orbitals.  The two defining physical properties follow from two
elementary properties of determinants:
\begin{itemize}
  \item interchanging two particles interchanges two \emph{columns}, which
        reverses the sign -- the wave function is antisymmetric, as
        Eq.~(\ref{eq:2-beta}) demands;
  \item placing two particles in the same orbital gives two identical
        \emph{rows}, and the determinant vanishes -- no two fermions may
        occupy the same single-particle state.
\end{itemize}
The Pauli principle is thus not an extra postulate bolted onto the formalism;
it is a property of determinants.

An equivalent form uses the antisymmetriser directly.  Writing the
unsymmetrised product, the \emph{Hartree function},
\begin{equation}
  \Phi_H(x_1,\dots,x_N;\alpha,\dots,\nu)
   = \psi_{\alpha}(x_1)\psi_{\beta}(x_2)\cdots\psi_{\nu}(x_N),
  \label{eq:2-hartreeproduct}
\end{equation}
we have
\begin{equation}
  \Phi = \frac{1}{\sqrt{N!}}\sum_{p}(-)^{p}\hat{P}\,\Phi_H
       = \sqrt{N!}\,\hat{A}\,\Phi_H .
  \label{eq:2-slaterasA}
\end{equation}
This is the form used in the derivations of Section~\ref{sec:expval}.

\begin{notebox}
\textbf{Link to Chapter~\ref{chap:linalg}.} The Hartree
function~(\ref{eq:2-hartreeproduct}) is precisely a tensor product,
$\bket{\alpha}\otimes\bket{\beta}\otimes\cdots\otimes\bket{\nu}$, of the kind
introduced in Section~\ref{sec:tensor}.  The Slater determinant is its
antisymmetric projection.  A tensor product of $N$ single-particle states is
the least entangled state there is, and its Schmidt rank across any cut is
one; antisymmetrisation introduces exactly the correlation demanded by
statistics and no more.  This is the precise sense in which a Slater
determinant is an \emph{uncorrelated} state, and it is made quantitative by
the occupation numbers of Section~\ref{sec:naturalorbitals}: for a single
Slater determinant they are exactly $1$ for the $N$ occupied orbitals and
exactly $0$ for all the rest.  Any deviation from those values is a measure of
genuine, dynamical correlation.
\end{notebox}

The numerical companion \texttt{slaterdeterminant.py} verifies both
statements.  Evaluating Eq.~(\ref{eq:SlatDet}) for three particles at
$(x_1,x_2,x_3)$ and at $(x_2,x_1,x_3)$ gives values that cancel to machine
zero, and putting two particles in the same orbital gives
$\Phi=-6\times10^{-18}$.  The eigenvalues of the density matrix
$\rho_{\lambda\mu}=\sum_i C^*_{i\lambda}C_{i\mu}$ come out as
$(1,1,1,0,0,0,0,0)$ to machine precision.

%----------------------------------------------------------------------
\section{Single-particle bases and unitary transformations}
\label{sec:basischange}

In practice one rarely knows the right orbitals in advance.  The standard
strategy is to expand them in a fixed, known basis $\{\phi_{\lambda}\}$ --
harmonic oscillator functions, hydrogen-like functions, plane waves,
Gaussians -- and to vary the expansion coefficients,
\begin{equation}
  \psi_p(x) = \sum_{\lambda} C_{p\lambda}\,\phi_{\lambda}(x).
  \label{eq:newbasis}
\end{equation}
The basis $\{\phi_{\lambda}\}$ is normally chosen as the eigenfunctions of the
one-body Hamiltonian, Eq.~(\ref{eq:2-spequation}), and is orthonormal.  In
Dirac notation Eq.~(\ref{eq:newbasis}) reads
\begin{equation}
  \bket{p} = \sum_{\lambda}C_{p\lambda}\bket{\lambda},
  \qquad
  \phi_{\lambda}(\bm{r}) = \braket{\bm{r}}{\lambda}.
  \label{eq:2-newbasisdirac}
\end{equation}

Two questions arise immediately.  Does the new basis remain orthonormal?  And
what happens to the Slater determinant?  Chapter~\ref{chap:linalg} has already
answered both.

\paragraph{Orthonormality is preserved if and only if $\bm{C}$ is unitary.}
This was proved in Section~\ref{sec:unitary}: for an orthonormal set
$\bm{v}_j^{T}\bm{v}_i=\delta_{ij}$ and a transformation
$\bm{w}_i=\bm{U}\bm{v}_i$,
\begin{equation}
  \bm{w}_j^{T}\bm{w}_i
   = \bm{v}_j^{T}\underbrace{\bm{U}^{T}\bm{U}}_{=\bm{I}}\bm{v}_i
   = \delta_{ij},
  \label{eq:2-orthopreserved}
\end{equation}
so $\braket{p}{q}=\delta_{pq}$ follows whenever $(C_{p\lambda})$ is unitary.
We shall use this constantly: it is what allows us to search over
single-particle bases without ever having to re-orthogonalise.

\paragraph{The Slater determinant picks up a phase.}
Substituting Eq.~(\ref{eq:newbasis}) into Eq.~(\ref{eq:SlatDet}), every entry
of the determinant becomes a sum $\sum_{\lambda}C_{p\lambda}\phi_{\lambda}(x_i)$,
that is, the matrix of orbitals is the product of the coefficient matrix
$\bm{C}$ with the matrix $\bm{\Phi}$ of basis functions.  By the product
rule~(\ref{eq:2-detproduct}),
\begin{equation}
  \frac{1}{\sqrt{N!}}
  \begin{vmatrix}
    \displaystyle\sum_\lambda C_{p\lambda}\phi_\lambda(x_1) & \cdots &
    \displaystyle\sum_\lambda C_{p\lambda}\phi_\lambda(x_N)\\
    \vdots & \ddots & \vdots\\
    \displaystyle\sum_\lambda C_{t\lambda}\phi_\lambda(x_1) & \cdots &
    \displaystyle\sum_\lambda C_{t\lambda}\phi_\lambda(x_N)
  \end{vmatrix}
  = \Det{C}\cdot\Det{\Phi} .
  \label{eq:2-slaterrotated}
\end{equation}
If $\bm{C}$ is unitary then $|\Det{C}|=1$, so the rotated Slater determinant
differs from the original only by an overall phase.  It describes the
\emph{same} physical state.

\begin{notebox}
\textbf{Many-body connection.} Equation~(\ref{eq:2-slaterrotated}) is the
algebraic foundation of a fact used everywhere in this book: one may perform
any unitary rotation among the occupied orbitals without changing the
many-body state, the density, or the total energy.  In Hartree-Fock theory
this is the freedom that makes canonical and localised orbitals equally valid
descriptions of the same determinant, and it is the content of Thouless'
theorem.  In configuration interaction and coupled-cluster theory it is why
the choice of reference orbitals is a matter of convenience rather than
principle.  The numerical check in \texttt{slaterdeterminant.py} makes it
concrete: a random orthogonal rotation of three occupied orbitals leaves the
energy unchanged to $2.7\times10^{-15}$, while a rotation that mixes occupied
and unoccupied orbitals changes it by more than a factor of two.
\end{notebox}

\paragraph{Truncation, and what happens when the basis is not orthogonal.}
A complete basis $\{\bket{\lambda}\}$ contains infinitely many states, and any
calculation must truncate the sum in Eq.~(\ref{eq:newbasis}).  This is the
first and often the largest approximation in a many-body calculation, and its
convergence has to be studied explicitly.

Furthermore, some of the most useful bases are not orthogonal at all.
Gaussian basis sets in quantum chemistry, oscillator states of different
oscillator length, and generator-coordinate states all have non-trivial
overlap matrices $S_{ij}=\braket{\phi_i}{\phi_j}$, and
Eq.~(\ref{eq:2-orthopreserved}) no longer applies.  The eigenvalue problem
then becomes the generalised one,
$\bm{H}\bm{c}=E\bm{S}\bm{c}$, treated in Section~\ref{sec:pseudoinverse}: one
diagonalises $\bm{S}$, discards the directions whose eigenvalues lie below a
threshold, and transforms with $\bm{X}=\bm{V}\bm{s}^{-1/2}$.  The threshold is
a physical choice, and Table~\ref{tab:canonical} shows how much
ill-conditioning it removes.

%----------------------------------------------------------------------
\section{Particle-hole notation and the Fermi level}
\label{sec:particlehole}

Before deriving the energy it is worth fixing the index conventions that will
be used for the rest of the book.  Given a reference Slater determinant --
usually the one obtained by filling the $N$ lowest single-particle levels --
the single-particle states divide into those that are occupied in the
reference and those that are not.  The dividing line is the \emph{Fermi
level}, denoted $F$.  We write
\begin{equation}
  \begin{array}{ll}
    i,j,k,l,\dots \le F & \text{occupied states, or \emph{hole} states},\\[2pt]
    a,b,c,d,\dots > F   & \text{unoccupied states, or \emph{particle} states},\\[2pt]
    p,q,r,s,\dots       & \text{general single-particle states.}
  \end{array}
  \label{eq:2-phnotation}
\end{equation}
The names come from the picture in which the reference determinant is treated
as a new vacuum: exciting a particle from an occupied state $i$ to an
unoccupied state $a$ creates a particle in $a$ and leaves a hole in $i$.
Chapter~\ref{chap:secondquant} develops this into the particle-hole formalism
of second quantisation, where it becomes the natural language for
perturbation theory, coupled-cluster theory and the shell model.  For now it
is simply a convention, but it is one worth adopting at once: every equation
from here on is easier to read with it.

%----------------------------------------------------------------------
\section{The variational principle}
\label{sec:variational}

Let $E_0$ be the exact ground-state energy of $\hat{H}$.  For any normalised
trial function $\Phi$,
\begin{equation}
  E_0 \le E[\Phi] = \int \Phi^{*}\hat{H}\,\Phi\,\ud{\bm{\tau}},
  \qquad
  \int \Phi^{*}\Phi\,\ud{\bm{\tau}} = 1,
  \label{eq:2-variational}
\end{equation}
with the shorthand $\ud{\bm{\tau}}=\ud{x_1}\ud{x_2}\cdots\ud{x_N}$.  The proof
is the spectral decomposition of Section~\ref{sec:spectral} and nothing more.
Expanding $\Phi$ in the exact eigenstates,
$\bket{\Phi}=\sum_k c_k\bket{\Psi_k}$ with $\hat{H}\bket{\Psi_k}=E_k
\bket{\Psi_k}$ and $\sum_k|c_k|^2=1$, gives
\begin{equation}
  E[\Phi] = \sum_k |c_k|^2 E_k
          \ \ge\ E_0\sum_k |c_k|^2 = E_0 ,
  \label{eq:2-variationalproof}
\end{equation}
with equality only when $\Phi$ is the ground state itself.

This one inequality is the engine of most of what follows.  Restricting the
trial function to a single Slater determinant and minimising gives
Hartree-Fock theory; allowing a linear combination of determinants gives
configuration interaction; parametrising it by a quantum circuit gives the
variational quantum eigensolver of Section~\ref{sec:expdim}.  It is also the
reason the Lanczos algorithm of Section~\ref{sec:lanczos} is so well suited to
many-body problems: the lowest Ritz value is the minimum of the Rayleigh
quotient over the Krylov space, and therefore an upper bound on $E_0$ that can
only improve as the space grows.

%----------------------------------------------------------------------
\section{Expectation values with a Slater determinant}
\label{sec:expval}

We now compute $E[\Phi]$ for the Slater determinant~(\ref{eq:SlatDet}).  The
calculation uses only Eqs.~(\ref{eq:2-commuteA}), (\ref{eq:Asq}) and the
orthonormality of the single-particle functions, and it is worth following in
full: the same manipulations recur throughout many-body theory, and second
quantisation exists largely to make them automatic.

\paragraph{The one-body contribution.}
Using the form~(\ref{eq:2-slaterasA}) of the Slater determinant,
\begin{equation}
  \int \Phi^{*}\hat{H}_0\Phi\,\ud{\bm{\tau}}
   = N!\int \Phi_H^{*}\hat{A}\hat{H}_0\hat{A}\,\Phi_H\,\ud{\bm{\tau}} .
  \label{eq:2-onebody1}
\end{equation}
Because $\hat{H}_0$ commutes with $\hat{A}$ we may move the left-hand
antisymmetriser through it, and because $\hat{A}^2=\hat{A}$ the two then
collapse into one,
\begin{equation}
  \int \Phi^{*}\hat{H}_0\Phi\,\ud{\bm{\tau}}
   = N!\int \Phi_H^{*}\hat{H}_0\hat{A}\,\Phi_H\,\ud{\bm{\tau}} .
  \label{eq:2-onebody2}
\end{equation}
Replacing $\hat{A}$ by its definition~(\ref{eq:2-antisym}), which cancels the
factor $N!$, and $\hat{H}_0$ by the sum of one-body operators,
\begin{equation}
  \int \Phi^{*}\hat{H}_0\Phi\,\ud{\bm{\tau}}
   = \sum_{i=1}^{N}\sum_{p}(-)^{p}
     \int \Phi_H^{*}\hat{h}_0(x_i)\hat{P}\,\Phi_H\,\ud{\bm{\tau}} .
  \label{eq:2-onebody3}
\end{equation}
Now the key step.  The integrand factorises into single-particle integrals,
and $\hat{h}_0(x_i)$ touches only one coordinate.  If $\hat{P}$ permutes any
two particles, at least one of the remaining factors is an overlap
$\braket{\mu}{\nu}$ between two \emph{different} orbitals, which vanishes.
Only the identity permutation survives,
\begin{equation}
  \int \Phi^{*}\hat{H}_0\Phi\,\ud{\bm{\tau}}
   = \sum_{i=1}^{N}\int \Phi_H^{*}\hat{h}_0(x_i)\Phi_H\,\ud{\bm{\tau}} ,
  \label{eq:2-onebody4}
\end{equation}
and integrating out the $N-1$ spectator coordinates, each of which
contributes unity, leaves
\begin{equation}
  \int \Phi^{*}\hat{H}_0\Phi\,\ud{\bm{\tau}}
   = \sum_{\mu=1}^{N}\int \psi_{\mu}^{*}(x)\,\hat{h}_0\,\psi_{\mu}(x)\ud{x} .
  \label{H1Expectation}
\end{equation}
Introducing the shorthand
\begin{equation}
  \langle \mu|\hat{h}_0|\mu\rangle
   = \int \psi_{\mu}^{*}(x)\hat{h}_0\psi_{\mu}(x)\ud{x},
  \label{eq:2-onebodyshorthand}
\end{equation}
this is simply
\begin{equation}
  \int \Phi^{*}\hat{H}_0\Phi\,\ud{\bm{\tau}}
   = \sum_{\mu=1}^{N}\langle\mu|\hat{h}_0|\mu\rangle .
  \label{H1Expectation1}
\end{equation}
The one-body energy is the sum of the diagonal matrix elements over the
occupied orbitals -- the result one would have guessed, now derived.

\paragraph{The two-body contribution.}
The same three steps apply,
\begin{equation}
  \int \Phi^{*}\hat{H}_I\Phi\,\ud{\bm{\tau}}
   = N!\int \Phi_H^{*}\hat{A}\hat{H}_I\hat{A}\,\Phi_H\,\ud{\bm{\tau}}
   = \sum_{i<j}^{N}\sum_{p}(-)^{p}
     \int \Phi_H^{*}\hat{v}(r_{ij})\hat{P}\,\Phi_H\,\ud{\bm{\tau}} ,
  \label{eq:2-twobody1}
\end{equation}
but the conclusion is different.  Because $\hat{v}(r_{ij})$ depends on
\emph{two} coordinates, the permutation that exchanges precisely those two
particles no longer produces a vanishing overlap.  Two permutations survive:
the identity, and the transposition $\hat{P}_{ij}$.  Hence
\begin{equation}
  \int \Phi^{*}\hat{H}_I\Phi\,\ud{\bm{\tau}}
   = \sum_{i<j}^{N}\int
     \Phi_H^{*}\,\hat{v}(r_{ij})\left(1-\hat{P}_{ij}\right)\Phi_H\,\ud{\bm{\tau}} .
  \label{eq:2-twobody2}
\end{equation}
Writing out the two terms and integrating over the spectator coordinates,
\begin{align}
  \int \Phi^{*}\hat{H}_I\Phi\,\ud{\bm{\tau}}
  = \frac{1}{2}\sum_{\mu=1}^{N}\sum_{\nu=1}^{N}
    &\left[\int \psi_{\mu}^{*}(x_i)\psi_{\nu}^{*}(x_j)\,
      \hat{v}(r_{ij})\,\psi_{\mu}(x_i)\psi_{\nu}(x_j)\,\ud{x_i}\ud{x_j}
     \right.\nonumber\\
  &\left.\ -\int \psi_{\mu}^{*}(x_i)\psi_{\nu}^{*}(x_j)\,
      \hat{v}(r_{ij})\,\psi_{\nu}(x_i)\psi_{\mu}(x_j)\,\ud{x_i}\ud{x_j}
    \right],
  \label{H2Expectation}
\end{align}
where the factor $\tfrac12$ compensates for the unrestricted double sum, which
counts every pair twice.  The first term is the \emph{direct} or
\emph{Hartree} term: it is the classical electrostatic energy of one charge
distribution in the field of another.  The second is the \emph{exchange} or
\emph{Fock} term, and it has no classical counterpart whatsoever; it exists
solely because the wave function is antisymmetric, and it is the sole trace
left in the energy by the transposition $\hat{P}_{ij}$ in
Eq.~(\ref{eq:2-twobody2}).

\paragraph{Notation for the matrix elements.}
Writing the single-particle functions as overlaps,
$\psi_{\alpha}(x)=\braket{x}{\alpha}$, we abbreviate the two integrals as
\begin{align}
  \langle\mu\nu|\hat{v}|\mu\nu\rangle
   &= \int \psi_{\mu}^{*}(x_i)\psi_{\nu}^{*}(x_j)\hat{v}(r_{ij})
      \psi_{\mu}(x_i)\psi_{\nu}(x_j)\,\ud{x_i}\ud{x_j},
  \label{eq:2-direct}\\
  \langle\mu\nu|\hat{v}|\nu\mu\rangle
   &= \int \psi_{\mu}^{*}(x_i)\psi_{\nu}^{*}(x_j)\hat{v}(r_{ij})
      \psi_{\nu}(x_i)\psi_{\mu}(x_j)\,\ud{x_i}\ud{x_j},
  \label{eq:2-exchange}
\end{align}
and define the \emph{antisymmetrised} matrix element
\begin{equation}
  \langle\mu\nu|\hat{v}|\mu\nu\rangle_{\mathrm{AS}}
   = \langle\mu\nu|\hat{v}|\mu\nu\rangle
   - \langle\mu\nu|\hat{v}|\nu\mu\rangle .
  \label{eq:2-antisymmatrix}
\end{equation}
Note that $\langle\mu\mu|\hat{v}|\mu\mu\rangle_{\mathrm{AS}}=0$: a particle
does not interact with itself.  The self-interaction that the Hartree term
would otherwise contain is cancelled exactly by the exchange term, which is
one of the reasons Hartree-Fock theory is better behaved than the Hartree
theory that preceded it.

%----------------------------------------------------------------------
\section{The energy as a functional of the coefficients}
\label{sec:energyfunctional}

Collecting Eqs.~(\ref{H1Expectation1}) and (\ref{H2Expectation}), the
expectation value of the Hamiltonian in a Slater determinant is
\begin{equation}
  E[\Phi]
   = \sum_{\mu=1}^{N}\langle\mu|\hat{h}_0|\mu\rangle
   + \frac{1}{2}\sum_{\mu=1}^{N}\sum_{\nu=1}^{N}
     \langle\mu\nu|\hat{v}|\mu\nu\rangle_{\mathrm{AS}} ,
  \label{eq:EHF}
\end{equation}
or, in the particle-hole notation of Section~\ref{sec:particlehole},
\begin{equation}
  E[\Phi] = \sum_{i=1}^{N}\langle i|\hat{h}_0|i\rangle
          + \frac{1}{2}\sum_{i,j=1}^{N}\langle ij|\hat{v}|ij\rangle_{\mathrm{AS}} .
  \label{FunctionalEPhi2}
\end{equation}
This is a functional of the occupied orbitals.  To turn it into something we
can minimise we insert the expansion~(\ref{eq:newbasis}) and let the
coefficients carry the variation,
\begin{equation}
  E[\Psi]
   = \sum_{i=1}^{N}\sum_{\alpha\beta}C^{*}_{i\alpha}C_{i\beta}
     \langle\alpha|\hat{h}_0|\beta\rangle
   + \frac{1}{2}\sum_{i,j=1}^{N}\sum_{\alpha\beta\gamma\delta}
     C^{*}_{i\alpha}C^{*}_{j\beta}C_{i\gamma}C_{j\delta}
     \langle\alpha\beta|\hat{v}|\gamma\delta\rangle_{\mathrm{AS}} .
  \label{eq:EHFnew}
\end{equation}
The Greek indices run over the full, fixed basis; the Latin indices run over
the $N$ occupied orbitals only.  Everything on the right-hand side except the
$C$'s is computed once and stored.

Minimising Eq.~(\ref{eq:EHFnew}) with respect to $C^{*}_{i\alpha}$ subject to
the orthonormality constraint yields the Hartree-Fock equations.  We defer
the derivation, but the structure of the answer can be read off immediately
and is worth stating now.  Define the \emph{density matrix} and the
\emph{Fock matrix}
\begin{equation}
  \rho_{\gamma\delta} = \sum_{i=1}^{N} C^{*}_{i\gamma}C_{i\delta},
  \qquad
  f_{\alpha\beta} = \langle\alpha|\hat{h}_0|\beta\rangle
    + \sum_{\gamma\delta}\rho_{\gamma\delta}
      \langle\alpha\gamma|\hat{v}|\beta\delta\rangle_{\mathrm{AS}} .
  \label{eq:2-fock}
\end{equation}
The condition for a stationary point is then
\begin{equation}
  \sum_{\beta} f_{\alpha\beta}\,C_{i\beta} = \varepsilon_i\,C_{i\alpha},
  \label{eq:2-hfequations}
\end{equation}
a Hermitian eigenvalue problem of exactly the kind solved in
Sections~\ref{sec:householder} and~\ref{sec:qlalgorithm} -- with the single
complication that the matrix $\bm{f}$ depends, through $\bm{\rho}$, on the
eigenvectors we are trying to find.  One therefore guesses, builds $\bm{f}$,
diagonalises, rebuilds and repeats until nothing changes.  This is the
\emph{self-consistent field} procedure, and it is an iteration in the sense of
Section~\ref{sec:iterative}, with a matrix diagonalisation in place of a
single sweep.

\begin{notebox}
\textbf{Summary: Hartree-Fock is Chapter~\ref{chap:linalg} applied to
fermions.} Expand the states in an orthonormal basis
(Section~\ref{sec:onb}); change basis with a unitary transformation
(Section~\ref{sec:unitary}), which leaves the determinant invariant up to a
phase by the product rule (Section~\ref{sec:determinants}); minimise a
quadratic form subject to a constraint, which produces a Hermitian eigenvalue
problem (Section~\ref{sec:spectral}); solve it by Householder reduction
followed by the QL algorithm (Sections~\ref{sec:householder}
and~\ref{sec:qlalgorithm}); and iterate to self-consistency.  Every step has
already been developed.
\end{notebox}

The minimal implementation in \texttt{slaterdeterminant.py} does exactly this,
using the eigenvalue solver written in Chapter~\ref{chap:linalg}.  For
spinless fermions in a one-dimensional trap with a softened Coulomb
repulsion, and a basis of eight harmonic oscillator orbitals, it converges in
nine iterations and lowers the energy below that of the naive determinant
built from the $N$ lowest oscillator states, as shown in
Table~\ref{tab:hfgain}.

\begin{table}[h]
\centering
\renewcommand{\arraystretch}{1.25}
\setlength{\tabcolsep}{10pt}
\begin{tabular}{rlll}
\toprule
$N$ & $E$ (lowest $N$ orbitals) & $E$ (Hartree-Fock) & \textbf{Gain}\\
\midrule
2 & $2.69018342$  & $2.66308878$  & $-2.7\times10^{-2}$\\
3 & $6.46801021$  & $6.39016133$  & $-7.8\times10^{-2}$\\
4 & $11.77084712$ & $11.62305385$ & $-1.5\times10^{-1}$\\
\bottomrule
\end{tabular}
\caption{Self-consistent Hartree-Fock energies for $N$ spinless fermions in a
one-dimensional harmonic trap with a softened Coulomb interaction, in a basis
of eight oscillator orbitals.  The middle column is the energy of the
determinant built from the $N$ lowest oscillator states; the Hartree-Fock
solution optimises the orbitals and always lies below it.}
\label{tab:hfgain}
\end{table}

For $N=2$ the exact ground state in the same truncated basis can be obtained
by diagonalising the Hamiltonian in the space of antisymmetric pairs, and the
difference
\begin{equation}
  E_{\mathrm{corr}} = E_{\mathrm{exact}} - E_{\mathrm{HF}}
  \label{eq:2-correlation}
\end{equation}
defines the \emph{correlation energy}.  Here it is $-3.86\times10^{-3}$: the
system is weakly correlated, and the largest amplitude in the exact state,
$0.9930$, sits on the Hartree-Fock configuration.  Recovering the remaining
$0.7\%$ of the amplitude is what every method in the rest of this book is
built to do.

%----------------------------------------------------------------------
\section{The price of the two-body matrix elements}
\label{sec:twobodycost}

Equation~(\ref{eq:EHFnew}) contains a sum over four basis indices, so a
calculation in a basis of $n$ single-particle states requires $n^4$ two-body
matrix elements.  For $n=100$ that is $10^8$ numbers, for $n=1000$ it is
$10^{12}$, and long before the many-body problem becomes hard the mere
storage of the interaction does.

Symmetries help.  For a real, local interaction the matrix elements obey
\begin{equation}
  \langle\alpha\beta|\hat{v}|\gamma\delta\rangle
   = \langle\beta\alpha|\hat{v}|\delta\gamma\rangle
   = \langle\gamma\delta|\hat{v}|\alpha\beta\rangle ,
  \label{eq:2-vsymmetry}
\end{equation}
which reduces the count by roughly a factor of eight, and rotational or
translational invariance can reduce it much further by making most elements
vanish.  But the scaling remains $\bigO(n^4)$.

The structural remedy is the one developed in
Section~\ref{sec:twobodyfact}.  Reading the pairs $(\alpha\gamma)$ and
$(\beta\delta)$ as composite indices turns the four-index array into an
ordinary matrix, which for a Coulomb-like interaction is positive
semidefinite with rapidly decaying eigenvalues, so that
\begin{equation}
  \langle\alpha\beta|\hat{v}|\gamma\delta\rangle
   = \sum_{\eta=0}^{M-1} L^{\eta}_{\alpha\gamma}L^{\eta}_{\beta\delta},
  \qquad M \ll n^2 .
  \label{eq:2-cholesky}
\end{equation}
The storage drops from $n^4$ to $Mn^2$ and the contractions in
Eq.~(\ref{eq:EHFnew}) can be performed on the factors rather than on the full
array.  For the eight-orbital basis used above, the $64\times64$ pair matrix
has numerical rank $15$, not $64$: less than a quarter of the nominal size.
This is the same effect that Table~\ref{tab:density} quantified, and it is
what makes large-basis coupled-cluster and perturbation-theory calculations
possible at all.

%----------------------------------------------------------------------
\section{Quantum computing: the gate model and its linear-algebra foundation}
\label{sec:qcgates}

A quantum computer is a well-controlled quantum many-body system, and the
abstract \emph{gate-based} model is defined entirely in the language of
Chapter~\ref{chap:linalg}.  Starting from $N$ qubits initialised in
$\bket{\Psi}^{(0)}=\bket{00\cdots0}$, the state evolves by successive
applications of unitary gates,
\[
  \bket{\Psi}^{(n+1)}=\hat{U}^{(n)}\bket{\Psi}^{(n)},\qquad
  \hat{U}^{(n)}=T\exp\!\left(-\imath\int_{t_n}^{t_{n+1}}\hat{H}(t)\,\ud{t}\right).
\]
A \emph{single-qubit gate} acts on qubit $a$ via a $2\times 2$ unitary:
\[
  \Psi_{i_1\cdots i_N}^{(n+1)}
  =\sum_{i'_a}U^{(n)}_{i_a i'_a}\,
   \Psi_{i_1\cdots i_{a-1}i'_a i_{a+1}\cdots i_N}^{(n)}.
\]
A \emph{two-qubit gate} acts on qubits $a$, $b$ via a $4\times 4$
unitary.  Standard examples are the Pauli gates $X,Y,Z$ of
Section~\ref{sec:pauli}, the Hadamard gate
\[
  H=\frac{1}{\sqrt{2}}\begin{pmatrix}1&1\\1&-1\end{pmatrix},
\]
the phase gate
$T=\begin{pmatrix}e^{\imath\pi/8}&0\\0&e^{-\imath\pi/8}\end{pmatrix}$,
and the two-qubit CNOT (controlled-NOT):
\[
  C_X=\begin{pmatrix}1&0&0&0\\0&1&0&0\\0&0&0&1\\0&0&1&0\end{pmatrix}.
\]
When qubit $a$ is \emph{measured}, the result $\alpha\in\{0,1\}$
occurs with probability
\[
  P_\alpha
  =\sum_{i_1\cdots\hat{\imath}_a\cdots i_N}
   \left|\Psi_{i_1\cdots\alpha\cdots i_N}^{(n)}\right|^2,
\]
after which the state collapses to the corresponding projected state -- the
projection operators of Section~\ref{sec:compbasis} acting for real.

\begin{notebox}
\textbf{In a nutshell.} Quantum computing is a restricted form of
linear algebra: matrix-vector multiplications with special unitary
matrices (gates) on exponentially large vectors (wave functions).
The challenge -- shared with classical many-body physics -- is that
only a tiny fraction of the $2^N$-dimensional space is physically
relevant; identifying that subspace is the core problem in both fields.
\end{notebox}

%----------------------------------------------------------------------
\section{The exponential-dimension problem and why it matters}
\label{sec:expdim}

The state of $N$ spin-$\frac12$ particles requires $2^N$ complex numbers to
describe, as the tensor-product construction of Section~\ref{sec:tensor}
already made clear.  For $N\approx 300$, $2^{300}\gg 10^{82}$, the estimated
number of atoms in the observable universe.  The same counting applies to the
fermionic problem: the number of Slater determinants that can be formed by
distributing $N$ particles over $n$ single-particle states is $\binom{n}{N}$,
which for $n=40$ and $N=20$ is already $1.4\times10^{11}$.

How do physicists actually compute?
\begin{enumerate}
  \item \textbf{Clever approximations}: Bethe ansatz, perturbation
        theory, mean-field approximations such as the Hartree-Fock theory
        derived above.
  \item \textbf{Computer simulations}: DFT, quantum Monte Carlo, FCI,
        coupled cluster, DMRG.
  \item \textbf{Restricting to tractable subsystems}: systems where a
        single-particle basis is a good approximation.
\end{enumerate}

Chapter~\ref{chap:linalg} supplied two of the reasons this is not hopeless.
The first is that we rarely want the whole spectrum: the Lanczos algorithm of
Section~\ref{sec:lanczos} extracts the lowest eigenvalues of a matrix of
dimension $10^9$ using only matrix-vector products, and the pairing-model
example there converged to machine precision in forty of them.  The second is
that physically relevant states are compressible: the Schmidt spectrum of a
ground state decays rapidly (Section~\ref{sec:schmidt}), and the
Eckart-Young theorem, Eq.~(\ref{eq:1-eckartyoung}), guarantees that
truncating it is the best approximation of its rank.  Density-matrix
renormalisation group and tensor-network methods are built on precisely that
observation.

For a quantum circuit of depth $D$ acting on $N$ qubits, the accessible
subspace has dimension $\bigO(DN)$, vastly smaller than $2^N$.  For $N=100$
and $D=10^6$, this subspace still has $10^{20}$ times fewer degrees of freedom
than the full Hilbert space.  The question is whether the \emph{physically
relevant} states lie in this reachable subspace -- a central open problem in
quantum many-body theory and quantum computing alike (see, e.g.,
\href{https://www.nature.com/articles/s42254-025-00813-9}{LaRocca~et~al.}).

The variational quantum eigensolver (VQE) tackles the problem by using a
parametrised quantum circuit as an ansatz for the ground state and minimising
$\langle\Psi(\bm{\theta})|\hat{H}|\Psi(\bm{\theta})\rangle$ classically over
the circuit parameters.  It is the quantum-computing form of the variational
principle~(\ref{eq:2-variational}) that underlies Hartree-Fock and
coupled-cluster theory: the bound $E(\bm{\theta})\ge E_0$ holds for every
choice of parameters, and the quality of the result is decided entirely by
whether the ansatz can reach the true ground state.

%=============================================================================
\section{Exercises}
%=============================================================================


\subsection*{Chapter 2: Many-body physics and quantum computing}
\addcontentsline{toc}{subsection}{Chapter 2: Many-body physics and quantum computing}

\begin{enumerate}

\item \textbf{Determinant product rule.}
  For $2\times 2$ matrices $\bm{B}$ and $\bm{C}$, define
  $a_{ij}=\sum_k b_{ik}c_{kj}$.  Show by direct calculation that
  $\det(\bm{A})=\det(\bm{B})\det(\bm{C})$.

\item \textbf{Slater determinant: $N=2$.}
  Write the Slater determinant for two electrons in orbitals
  $\psi_\alpha$ and $\psi_\beta$.
  (a) Show it is antisymmetric under particle exchange.
  (b) Show it vanishes if $\psi_\alpha=\psi_\beta$.

\item \textbf{Antisymmetriser.}
  Verify for $N=2$ that $\hat{A}^2=\hat{A}$ and that
  $[\hat{H}_0,\hat{A}]=0$ when $\hat{H}_0=\hat{h}_0(x_1)+\hat{h}_0(x_2)$.

\item \textbf{Basis change of the Slater determinant.}
  With the expansion $\psi_p=\sum_\lambda C_{p\lambda}\phi_\lambda$
  and a $2\times 2$ unitary $\bm{C}$, show that the two-particle
  Slater determinant in the new basis equals $\det(\bm{C})\det(\bm{\Phi})$
  and that $|\det(\bm{C})|=1$ leaves the physics unchanged.

\item \textbf{Hartree-Fock energy.}
  For a two-electron system with orbitals $\psi_\alpha$ and
  $\psi_\beta$, write out the HF energy functional~\eqref{eq:EHF}
  explicitly, identifying the kinetic, direct (Hartree), and exchange
  (Fock) contributions.

\item \textbf{Variational principle.}
  Show that $E[\Phi]\geq E_0$ for any normalised trial state $\Phi$,
  where $E_0$ is the exact ground state energy.  Use the spectral
  decomposition of $\hat{H}$ in its eigenbasis.

\item \textbf{Exponential dimension.}
  How many complex numbers are needed to store the full $N$-particle
  wave function for $N=10$, $20$, $50$?  A modern GPU can perform
  roughly $10^{14}$ floating-point operations per second.  Estimate
  how long a single matrix-vector multiplication in the $N=50$ space
  would take.

\item \textbf{One-qubit gates numerically.}
  Write a short Python function that represents $\bket{0}$ and
  $\bket{1}$ as NumPy arrays and applies each Pauli matrix,
  the Hadamard gate $H$, and the phase gate $T$ in turn.
  Verify the action of $X$ (spin flip), $Z$ (phase flip) and
  $H\bket{0}=(\bket{0}+\bket{1})/\sqrt{2}$ (creation of a
  superposition state).

\item \textbf{Pauli gate as a quantum NOT.}
  Explain why $\sigma_x$ is the quantum analogue of the classical NOT
  gate.  What is the difference between the classical and quantum
  cases when the input is a superposition
  $\alpha\bket{0}+\beta\bket{1}$?

\item \textbf{VQE and the variational principle.}
  The variational quantum eigensolver minimises
  $E(\bm{\theta})=\langle\Psi(\bm{\theta})|\hat{H}|\Psi(\bm{\theta})\rangle$
  over circuit parameters $\bm{\theta}$.
  (a) Explain why $E(\bm{\theta})\geq E_0$ for any $\bm{\theta}$.
  (b) Describe the connection to the Hartree-Fock variational principle.
  (c) What determines the quality of the approximation in VQE
      compared to FCI?


\item \textbf{The antisymmetriser as a projector.}
  Using Eq.~(\ref{eq:2-antisym}), verify for $N=2$ that
  $\hat{A}^2=\hat{A}$ and that $[\hat{H}_0,\hat{A}]=0$ when
  $\hat{H}_0=\hat{h}_0(x_1)+\hat{h}_0(x_2)$.  Explain why
  $\hat{A}$ is a projection operator in the sense of
  Section~\ref{sec:compbasis}, and identify the subspace it projects onto.

\item \textbf{Why only two permutations survive.}
  In the derivation of Eq.~(\ref{eq:2-twobody2}) it was argued that only the
  identity and the transposition $\hat{P}_{ij}$ contribute to the two-body
  expectation value.  Write out the case $N=3$ explicitly, listing all six
  permutations, and show that four of them give a vanishing integral.
  Then explain why the corresponding argument for the one-body operator
  leaves only one surviving permutation.

\item \textbf{Self-interaction.}
  Show from Eq.~(\ref{eq:2-antisymmatrix}) that
  $\langle\mu\mu|\hat{v}|\mu\mu\rangle_{\mathrm{AS}}=0$, and explain in words
  what would go wrong in Eq.~(\ref{eq:EHF}) if the exchange term were omitted.

\item \textbf{Occupation numbers of a Slater determinant.}
  For the determinant~(\ref{eq:SlatDet}) the density matrix is
  $\rho_{\lambda\mu}=\sum_{i}C^{*}_{i\lambda}C_{i\mu}$.
  (a) Show that $\bm{\rho}^2=\bm{\rho}$ when the coefficients are orthonormal.
  (b) Conclude that the eigenvalues -- the natural occupation numbers of
      Section~\ref{sec:naturalorbitals} -- can only be $0$ or $1$.
  (c) What does it mean physically when a computed occupation number comes
      out as $0.93$?

\item \textbf{Cost of a determinant.}
  Reproduce Table~\ref{tab:detcost}.  For $n=20$, estimate how long the
  evaluation of a single determinant would take from
  Eq.~(\ref{eq:2-detdef}) on a machine performing $10^{12}$ operations per
  second, and compare with the LU route of Section~\ref{sec:lu}.

\item \textbf{Two-body storage.}
  How many two-body matrix elements $\langle\alpha\beta|\hat{v}|\gamma\delta
  \rangle$ are there for $n=50$, $n=200$ and $n=1000$ single-particle states?
  How much memory does each case require in double precision?  Using
  Eq.~(\ref{eq:2-cholesky}) with $M=5n$, how much is saved?

\item \textbf{Hartree-Fock numerically.}
  Using \texttt{slaterdeterminant.py}, reproduce Table~\ref{tab:hfgain}.
  Then vary the interaction strength and plot the correlation
  energy~(\ref{eq:2-correlation}) and the largest natural occupation number
  against it.  At what point would you no longer trust a single-determinant
  description?

\end{enumerate}

